% GENERATED FILE. Edit canonical lessons, then run npm run generate:books.
% canonical-source-hash: fnv1a64:e6f2309077723a4c
% canonical-lessons: FR-C328-planche-a-decouper, FR-C328-eplucher, FR-C328-hacher, FR-C328-emincer, FR-C328-raper

\chapter{Couper les légumes}
\label{ch:fr-couper-les-legumes}
\hlchaptermodality{328}

\begin{chapteropening}
\textbf{By the end of this chapter:} \emph{I can describe getting vegetables ready: the chopping board, peeling, chopping, slicing thinly and grating.}

Complete the last lesson of chapter 328: I can describe getting vegetables ready: the chopping board, peeling, chopping, slicing thinly and grating.
\end{chapteropening}

\section[la planche à découper]{la planche à découper — the chopping board}

\label{lesson:FR-C328-planche-a-decouper}



\begin{quote}
Before the new one: say the French for the spoonful, then the French for to preheat.
\end{quote}

\subsection*{You'll want to know: la planche à découper}
\textbf{la planche à découper} — \textquotedblleft{}the chopping board\textquotedblright{}.

\begin{quote}
Les légumes sont sur la planche à découper.
\end{quote}

\begin{quote}
\emph{The vegetables are on the chopping board.}
\end{quote}

\emph{Une planche} alone is a plank or a board.

\begin{cousinweb}
\emph{planche} is inherited from Late Latin \emph{planca} \textquotedblleft{}a board\textquotedblright{} (EtymDB); \emph{découper} is built on \emph{couper} \textquotedblleft{}to cut\textquotedblright{}, from \emph{coup} \textquotedblleft{}a blow\textquotedblright{}.
\end{cousinweb}

\begin{grammarlens}[title={What you've built}]
One more word for telling, explaining and describing at B1.
\end{grammarlens}

\subsection*{Your turn}
\begin{itemize}
\raggedright
  \item \emph{Say it:} \emph{la planche à découper}
  \item \emph{Say it:} \emph{la planche à découper}, once more
  \item \emph{Say it:} a sentence of your own with \emph{la planche à découper}
  \item \emph{Recall:} say the French for the spoonful, then the French for to preheat, then say \emph{la planche à découper} again
\end{itemize}

\begin{tcolorbox}[breakable,colback=teal!4,colframe=teal!35!black,title={Before you move on}]
What does la planche à découper mean? (\textquotedblleft{}The chopping board\textquotedblright{}.) Say it once more.
\end{tcolorbox}
\section[éplucher]{éplucher — to peel}

\label{lesson:FR-C328-eplucher}



\begin{quote}
Before the new one: say the French for to preheat, then the French for the chopping board.
\end{quote}

\subsection*{You'll want to know: éplucher}
\textbf{éplucher} — \textquotedblleft{}to peel\textquotedblright{}.

\begin{quote}
Mon père épluche les pommes de terre.
\end{quote}

\begin{quote}
\emph{My father is peeling the potatoes.}
\end{quote}

Figurative: \emph{éplucher un texte} — to go through it in detail.

\begin{cousinweb}
\emph{éplucher} is built on Old French \emph{peluchier} \textquotedblleft{}to pick at\textquotedblright{}, from Vulgar Latin \emph{piluccāre} (EtymDB), a verb usually traced to \emph{pilus} \textquotedblleft{}a hair\textquotedblright{} (Lewis \& Short).
\end{cousinweb}

\begin{grammarlens}[title={What you've built}]
One more word for telling, explaining and describing at B1.
\end{grammarlens}

\subsection*{Your turn}
\begin{itemize}
\raggedright
  \item \emph{Say it:} \emph{éplucher}
  \item \emph{Say it:} \emph{éplucher}, once more
  \item \emph{Say it:} a sentence of your own with \emph{éplucher}
  \item \emph{Recall:} say the French for to preheat, then the French for the chopping board, then say \emph{éplucher} again
\end{itemize}

\begin{tcolorbox}[breakable,colback=teal!4,colframe=teal!35!black,title={Before you move on}]
What does éplucher mean? (\textquotedblleft{}To peel\textquotedblright{}.) Say it once more.
\end{tcolorbox}
\section[hacher]{hacher — to chop, to mince}

\label{lesson:FR-C328-hacher}



\begin{quote}
Before the new one: say the French for the chopping board, then the French for to peel.
\end{quote}

\subsection*{You'll want to know: hacher}
\textbf{hacher} — \textquotedblleft{}to chop, to mince\textquotedblright{}.

\begin{quote}
Elle hache les oignons très fin.
\end{quote}

\begin{quote}
\emph{She chops the onions very finely.}
\end{quote}

\emph{La viande hachée} — minced meat; \emph{un steak haché} — a burger patty.

\begin{cousinweb}
\emph{hacher} is built on \emph{hache} \textquotedblleft{}axe\textquotedblright{}, which comes from a Frankish word \emph{hapja} (EtymDB).
\end{cousinweb}

\begin{grammarlens}[title={What you've built}]
One more word for telling, explaining and describing at B1.
\end{grammarlens}

\subsection*{Your turn}
\begin{itemize}
\raggedright
  \item \emph{Say it:} \emph{hacher}
  \item \emph{Say it:} \emph{hacher}, once more
  \item \emph{Say it:} a sentence of your own with \emph{hacher}
  \item \emph{Recall:} say the French for the chopping board, then the French for to peel, then say \emph{hacher} again
\end{itemize}

\begin{tcolorbox}[breakable,colback=teal!4,colframe=teal!35!black,title={Before you move on}]
What does hacher mean? (\textquotedblleft{}To chop, to mince\textquotedblright{}.) Say it once more.
\end{tcolorbox}
\section[émincer]{émincer — to slice thinly}

\label{lesson:FR-C328-emincer}



\begin{quote}
Before the new one: say the French for to peel, then the French for to chop.
\end{quote}

\subsection*{You'll want to know: émincer}
\textbf{émincer} — \textquotedblleft{}to slice thinly\textquotedblright{}.

\begin{quote}
Le cuisinier émince les poivrons.
\end{quote}

\begin{quote}
\emph{The cook slices the peppers thinly.}
\end{quote}

A cooking word; \emph{couper en fines tranches} says the same.

\begin{cousinweb}
\emph{émincer} is built on \emph{mince} \textquotedblleft{}thin\textquotedblright{}, from Old French \emph{mincier} \textquotedblleft{}to cut fine\textquotedblright{}, usually traced to Vulgar Latin \emph{minūtiāre}, built on \emph{minūtus} \textquotedblleft{}small\textquotedblright{} (Lewis \& Short).
\end{cousinweb}

\begin{grammarlens}[title={What you've built}]
One more word for telling, explaining and describing at B1.
\end{grammarlens}

\subsection*{Your turn}
\begin{itemize}
\raggedright
  \item \emph{Say it:} \emph{émincer}
  \item \emph{Say it:} \emph{émincer}, once more
  \item \emph{Say it:} a sentence of your own with \emph{émincer}
  \item \emph{Recall:} say the French for to peel, then the French for to chop, then say \emph{émincer} again
\end{itemize}

\begin{tcolorbox}[breakable,colback=teal!4,colframe=teal!35!black,title={Before you move on}]
What does émincer mean? (\textquotedblleft{}To slice thinly\textquotedblright{}.) Say it once more.
\end{tcolorbox}
\section[râper]{râper — to grate}

\label{lesson:FR-C328-raper}



\begin{quote}
Before the new one: say the French for to chop, then the French for to slice thinly.
\end{quote}

\subsection*{You'll want to know: râper}
\textbf{râper} — \textquotedblleft{}to grate\textquotedblright{}.

\begin{quote}
Il râpe du fromage sur les pâtes.
\end{quote}

\begin{quote}
\emph{He grates cheese over the pasta.}
\end{quote}

\emph{Le fromage râpé} — grated cheese; the tool is \emph{la râpe}.

\begin{cousinweb}
\emph{râper} comes from a Germanic verb meaning \textquotedblleft{}to scrape together\textquotedblright{} (EtymDB lists Old High German \emph{raspōn} beside it).
\end{cousinweb}

\begin{grammarlens}[title={What you've built}]
One more word for telling, explaining and describing at B1.
\end{grammarlens}

\subsection*{Your turn}
\begin{itemize}
\raggedright
  \item \emph{Say it:} \emph{râper}
  \item \emph{Say it:} \emph{râper}, once more
  \item \emph{Say it:} a sentence of your own with \emph{râper}
  \item \emph{Recall:} say the French for to chop, then the French for to slice thinly, then say \emph{râper} again
\end{itemize}

\begin{tcolorbox}[breakable,colback=teal!4,colframe=teal!35!black,title={Before you move on}]
What does râper mean? (\textquotedblleft{}To grate\textquotedblright{}.) Say it once more.
\end{tcolorbox}
